\customchapter{TRABAJOS FUTUROS}

\begin{itemize}
\item \textbf{Extensión a otras tipologías patrimoniales.} Validar el modelo
híbrido sobre edificaciones de tipología distinta a la de catedral colonial
(conventos, huacas, arquitectura republicana) para medir la pérdida de
rendimiento fuera de la generalización declarada.

\item \textbf{Segmentación de instancias.} Pasar del etiquetado por punto a la
individualización de elementos (columna~1, columna~2), requisito para la
generación automática del modelo BIM.

\item \textbf{Ornamento fino.} Explorar esquemas de resolución adaptativa que
levanten la limitación del vóxel de 5\,cm sobre elementos decorativos por debajo
de esa escala.

\item \textbf{Preentrenamiento específico de dominio.} Evaluar si un
preentrenamiento auto-supervisado sobre un corpus de levantamientos
patrimoniales mejora las representaciones respecto de los pesos públicos de
Sonata, entrenados sobre escenas interiores genéricas.

\item \textbf{Vocabulario abierto.} Integrar modelos de lenguaje-visión para
nombrar automáticamente las clases descubiertas, en lugar de dejarlas como
agrupaciones sin etiqueta semántica.
\end{itemize}
